\chapter{Du signal à la prévision}\label{ch:rs:from-signal-to-forecast}

Le score composite d'une action vaut $+2$. Son rendement espéré le mois suivant vaut-il deux points de base ou deux cents~? Le composite du
chapitre~14 classe les actions~; un optimiseur de portefeuille, un budget de risque et un trader ont besoin de rendements espérés exprimés en unités de rendement, par action
et par horizon, et un gérant doit savoir quel ratio d'information les prévisions peuvent produire. Ce chapitre transforme les scores en
rendements espérés (calibration), énonce la loi qui relie la qualité d'une prévision au ratio d'information d'un portefeuille, et mesure sur une
prévision implantée où la promesse de la loi se perd. La réponse à la question d'ouverture, pour un score dont l'IC vaut 0.05 sur une action dont
la volatilité résiduelle mensuelle vaut 7.7~\%, soit 77 points de base.

\section{Calibrer un score}

\begin{definition}[Calibration des prévisions, règle de mise à l'échelle de l'alpha]\label{def:rs:from-signal-to-forecast:calibration}
La \emph{calibration des prévisions}\index{calibration des prévisions} est l'application qui associe un score à un rendement espéré tel que, parmi les titres
associés à une même prévision, le rendement réalisé moyen est égal à celle-ci. La \emph{règle de mise à l'échelle de l'alpha}\index{règle de mise à l'échelle de l'alpha} fixe le
\omterm{def:rs:anatomy-of-a-predictor:residual}{rendement résiduel} espéré d'un titre à $\alpha_i = \mathrm{IC}\cdot\omega_i\cdot z_i$, avec $z_i$ le score standardisé et $\omega_i$ le
volatilité résiduelle du titre sur l'horizon.
\end{definition}

\begin{proposition}[La règle de mise à l'échelle est une régression]\label{prop:rs:from-signal-to-forecast:scaling}
Si le score standardisé $z$ et le \omterm{def:rs:anatomy-of-a-predictor:residual}{rendement résiduel} normalisé par la volatilité $r/\omega$ sont conjointement gaussiens, de corrélation $\mathrm{IC}$ et
de variances unitaires, le rendement espéré conditionnellement au score vaut $\mathrm{E}[r \mid z] = \mathrm{IC}\cdot\omega\cdot z$.
\end{proposition}

\begin{proof}
Pour des variables standardisées conjointement gaussiennes, $\mathrm{E}[r/\omega \mid z] = \mathrm{IC}\cdot z$~; il suffit de multiplier par $\omega$.
\end{proof}

La règle rend trois hypothèses visibles~: que le score est standardisé, que l'IC est le même pour tous les titres (de sorte que les titres plus risqués
reçoivent des prévisions proportionnellement plus grandes), et que la relation est linéaire. Chacune peut être vérifiée. La prévision implantée du chapitre
(\texttt{rs\_forecast}) ajoute aux rendements mensuels ajustés du marché des 704 titres que \texttt{firm.synthmkt} recense pendant dix ans un
rendement espéré vrai de $0.05\,\omega_i s_i$, avec $s_i$ un score gaussien standard tiré à nouveau chaque mois et $\omega_i$ la volatilité résiduelle du titre
(médiane de 7.7~\% par mois). Estimé sur les cinq premières années, l'IC vaut 0.055. Sur les cinq dernières, la régression de Mincer--Zarnowitz des rendements réalisés
sur les prévisions mises à l'échelle (livre~4, chapitre~18) a une pente de 0.92, proche de 1 pour une prévision calibrée, et un $R^2$ de
0.24~\%~: les prévisions calibrées n'expliquent presque rien de la variance des rendements, et c'est ce que signifie un IC de 0.05. La même régression sur le
score brut a une pente de 0.0039, soit l'IC multiplié par une volatilité typique~: le score n'a jamais été exprimé en unités de rendement.

\begin{definition}[Courbe de calibration, régression isotone]\label{def:rs:from-signal-to-forecast:curve}
Une \emph{courbe de calibration}\index{courbe de calibration} représente le rendement réalisé moyen des titres regroupés par prévision (en classes) en fonction de la
prévision moyenne. La \emph{régression isotone}\index{régression isotone} ajuste la meilleure fonction non décroissante du score aux rendements réalisés par les
moindres carrés, en regroupant les groupes voisins qui violent l'ordre (Ayer, Brunk, Ewing, Reid et Silverman, 1955).
\end{definition}

Regroupés par décile de score, les rendements réalisés normalisés par la volatilité suivent la droite de la règle de mise à l'échelle, le bruit de 4\,224 observations par décile
(la \cref{fig:rs:from-signal-to-forecast:calibration}). La \omterm{def:rs:from-signal-to-forecast:curve}{régression isotone} donne une calibration sans l'hypothèse linéaire et la paye en
paliers~: des plateaux là où les données ne pouvaient pas ordonner des scores voisins, et des valeurs plus grandes dans les queues, ajustées sur peu de titres.
Entre les deux, une calibration linéaire à pente estimée est le choix par défaut~; la courbe sert de contrôle.

\begin{omfigure}
\begin{tikzpicture}
\begin{axis}[width=10cm,height=6cm,xlabel={score $z$},ylabel={rendement réalisé / volatilité résiduelle},tick label style={font=\small},
  label style={font=\small},xmin=-2.6,xmax=2.6,ymin=-0.22,ymax=0.14,grid=major,grid style={gray!25},
  yticklabel style={/pgf/number format/fixed,/pgf/number format/precision=2},scaled y ticks=false,legend pos=north west,
  legend style={font=\footnotesize},table/col sep=comma]
\addplot[omThm,thick,dashed] table[x=score,y=predicted]{figdata/research/15-from-signal-to-forecast/calibration.csv};
\addplot[only marks,mark=*,omDef,mark size=2pt] table[x=score,y=realised]{figdata/research/15-from-signal-to-forecast/calibration.csv};
\addplot[black!60,thick,const plot] table[x=score,y=fit]{figdata/research/15-from-signal-to-forecast/isotonic.csv};
\legend{règle de mise à l'échelle $\widehat{\mathrm{IC}}\cdot z$,déciles de score (réalisés),régression isotone}
\end{axis}
\end{tikzpicture}

\omcaption{Calibration du score implanté sur les cinq dernières années~: rendement réalisé moyen rapporté à la volatilité résiduelle, par décile de score, comparé à la
règle de mise à l'échelle avec l'IC estimé sur les cinq premières années (0.055), et l'ajustement isotone. Données~: \texttt{rs\_forecast} sur
\texttt{firm.synthmkt}.}\label{fig:rs:from-signal-to-forecast:calibration}
\end{omfigure}

\section{Mise à l'échelle entre actifs et horizons}

La règle met la prévision à l'échelle de la volatilité de chaque titre, de sorte qu'une prévision n'est comparable entre actifs qu'à travers son IC. Un signal qui est
aussi informatif sur une action de services publics que sur une action de biotechnologie prévoit un rendement plus élevé pour l'action de biotechnologie, parce que ses rendements sont
plus élevés~; une prévision donnant le même rendement espéré aux deux impliquerait un IC bien plus élevé pour l'action de services publics. Selon les horizons, la mise à l'échelle suit
la \omterm{def:rs:half-life-decay-and-stability:decay}{courbe de décroissance de l'IC} du chapitre~13~: le rendement espéré sur $h$ périodes vaut $\omega\,z$ fois la somme des IC à une période jusqu'à $h$ (pour
volatilité par période $\omega$), et non la prévision d'une période multipliée par $h$. Un signal dont l'information ne dure qu'un jour prévoit le même rendement espéré
sur un mois que sur un jour. Entre classes d'actifs, la standardisation doit se faire dans chaque classe, et l'IC être estimé par classe~; un
mélange d'une prévision actions et d'une prévision obligataire ajoute des rendements espérés, jamais des scores.

\section{La loi fondamentale de la gestion active}

\begin{definition}[Largeur, loi fondamentale de la gestion active]\label{def:rs:from-signal-to-forecast:law}
La \emph{largeur}\index{largeur} (breadth) d'une stratégie est le nombre de prévisions indépendantes sur lesquelles elle agit par an. La \emph{loi fondamentale de la
gestion active}\index{loi fondamentale de la gestion active} (Grinold, 1989) énonce que le ratio d'information du portefeuille optimal construit à partir de
prévisions de \omterm{def:rs:anatomy-of-a-predictor:ic}{coefficient d'information} $\mathrm{IC}$ et de largeur $\mathrm{BR}$ vaut $\mathrm{IR} \approx \mathrm{IC}\sqrt{\mathrm{BR}}$.
\end{definition}

Avec 704 titres prévus chaque mois et un IC de 0.05, la loi promet $0.05\sqrt{12 \times 704} = 4.60$. Le portefeuille qui exploite la prévision implantée
détient chaque titre en proportion de $\alpha_i/\omega_i^2$ (les pondérations moyenne-variance pour des résidus non corrélés), rééquilibré chaque mois. Lorsque
les résidus sont remplacés par un bruit gaussien de mêmes volatilités, indépendant d'un titre à l'autre, le ratio d'information réalisé du portefeuille sur dix
ans est de 4.92, dans la marge d'erreur d'échantillonnage de la loi (environ 0.45 pour dix ans de rendements mensuels). La loi s'applique là où ses hypothèses sont satisfaites.

\section{Où la loi cesse de tenir}

\begin{definition}[Largeur effective, coefficient de transfert]\label{def:rs:from-signal-to-forecast:breadth}
La \emph{largeur effective}\index{largeur effective} d'un ensemble de prévisions est le nombre de prévisions indépendantes qui donneraient le même
ratio d'information~; $n$ paris de corrélation moyenne deux à deux $\rho$ valent environ $n/(1 + (n - 1)\rho)$. Le \emph{coefficient de
transfert}\index{coefficient de transfert} d'un portefeuille est la corrélation transversale entre ses pondérations actives ajustées du risque $w_i\omega_i$
et les prévisions ajustées du risque $\alpha_i/\omega_i$.
\end{definition}

\begin{proposition}[La loi généralisée]\label{prop:rs:from-signal-to-forecast:generalised}
Sous les hypothèses de la loi, un portefeuille dont les pondérations actives ajustées du risque ont une corrélation $\mathrm{TC}$ avec les prévisions ajustées du risque
a un ratio d'information espéré égal à $\mathrm{TC}\cdot\mathrm{IC}\sqrt{\mathrm{BR}}$.
\end{proposition}

\begin{proof}
Posons $x_i = w_i\omega_i$ et $u_i = r_i/\omega_i$, de sorte que le rendement actif vaut $\sum_i x_iu_i$ avec $\mathrm{E}[u_i] = \mathrm{IC}\,z_i$ et un bruit indépendant
de variance unitaire. Pour un $x$ fixé avec $\sum x_i^2 = 1$, le rendement actif espéré vaut $\mathrm{IC}\sum x_iz_i = \mathrm{IC}\sqrt n\,\mathrm{TC}$ (la
corrélation de $x$ avec $z$, les deux normes étant fixées) et le risque vaut 1~; sur $\mathrm{BR}/n$ périodes indépendantes par an, le ratio est multiplié par
$\sqrt{\mathrm{BR}/n}$, soit $\mathrm{TC}\cdot\mathrm{IC}\sqrt{\mathrm{BR}}$.
\end{proof}

Clarke, de Silva et Thorley (2002) ont établi cette version généralisée pour les portefeuilles contraints, avec une décomposition ex post de la performance en
la réussite des prévisions et le bruit des contraintes. Sur le marché étudié dans ce chapitre, le portefeuille perd sa promesse en trois étapes
(la \cref{fig:rs:from-signal-to-forecast:shortfall})~:

\textbf{L'IC est bruité.} Qian et Hua ont reformulé la loi avec un IC aléatoire~: le ratio d'information est la moyenne de l'IC sur son écart
type (tel que résumé par Zhang, Wang et Cao, 2021). Avec les résidus propres au marché (facteurs sectoriels et de style, regroupement de
volatilité), l'IC mensuel a une moyenne de 0.0505 et un écart type de 0.0424, contre 0.0331 avec des résidus gaussiens~; le ratio, annualisé, vaut 4.13, et le
ratio d'information réalisé du portefeuille vaut 4.26. Les facteurs font bouger ensemble des paris voisins, et les 704 titres valent
$(\mu_{\mathrm{IC}}/\sigma_{\mathrm{IC}})^2/\mathrm{IC}^2 = 569$ prévisions indépendantes par mois.

\textbf{Le portefeuille ne peut pas suivre les prévisions.} Un portefeuille long only construit autour d'un indice de référence à pondération égale ne peut pas détenir un titre à un poids négatif. Si on lui impose un
tracking error mensuel de 4~\%, il tronque à zéro 47~\% des positions visées, et son \omterm{def:rs:from-signal-to-forecast:breadth}{coefficient de transfert} tombe à 0.86. La loi généralisée
prédit alors $0.857 \times 4.13 = 3.54$.

\textbf{La contrainte ajoute du bruit.} Le ratio d'information réalisé du portefeuille est 3.00~: la part du portefeuille tronqué qui n'est pas proportionnelle aux
prévisions est non corrélée avec elles et ajoute du risque sans rendement. De 4.60 à 3.00, deux tiers de la promesse subsistent.

\begin{omfigure}
\begin{tikzpicture}
\begin{axis}[width=13.5cm,height=5.8cm,ybar,bar width=18pt,ylabel={ratio d'information (annuel)},tick label style={font=\small},
  label style={font=\small},xtick={0,1,2,3,4,5},
  xticklabels={{loi},{portefeuille\\gaussien},{IC moyen /\\écart type},{portefeuille\\marché},{$\times$ TC,\\long-only},{portefeuille,\\long-only}},
  xticklabel style={align=center,font=\footnotesize},ymin=0,ymax=5.5,enlarge x limits=0.1,grid=major,grid style={gray!25},
  nodes near coords,nodes near coords style={font=\footnotesize,/pgf/number format/fixed,/pgf/number format/fixed zerofill,
  /pgf/number format/precision=2},table/col sep=comma]
\addplot[fill=omDef!55,draw=omDef] table[x=k,y=ir]{figdata/research/15-from-signal-to-forecast/shortfall.csv};
\end{axis}
\end{tikzpicture}

\omcaption{Ratio d'information de la prévision artificielle (IC 0.05, 704 titres, mensuel)~: la promesse de la loi, celle du portefeuille avec des résidus gaussiens, le
rapport de la moyenne de l'IC à son écart type, celui du portefeuille avec les résidus du marché, la loi généralisée avec le \omterm{def:rs:from-signal-to-forecast:breadth}{coefficient de transfert} long-only, et
le portefeuille long-only. Données~: \texttt{rs\_forecast.law}.}\label{fig:rs:from-signal-to-forecast:shortfall}
\end{omfigure}

Le tracking error est le levier~: le \omterm{def:rs:from-signal-to-forecast:breadth}{coefficient de transfert} passe de 0.98 pour un tracking error mensuel demandé de 0.25~\% à 0.86 pour 4~\%, et le ratio d'information
passe de 3.91 à 3.00, tandis que la part des positions tronquées à zéro monte de 13~\% à 47~\%.

\begin{center}
\small
\begin{tabular}{@{}lccccc@{}}
\toprule
tracking error demandé (par mois) & 0.25~\% & 0.5~\% & 1~\% & 2~\% & 4~\%\\
\midrule
positions tronquées à zéro & 13~\% & 28~\% & 38~\% & 44~\% & 47~\%\\
\omterm{def:rs:from-signal-to-forecast:breadth}{coefficient de transfert} & 0.98 & 0.93 & 0.89 & 0.87 & 0.86\\
ratio d'information & 3.91 & 3.54 & 3.25 & 3.09 & 3.00\\
\bottomrule
\end{tabular}
\end{center}

\section{Évaluation des prévisions}

Une prévision s'évalue d'abord comme prévision avant d'être évaluée comme stratégie~: son IC (moyenne, dispersion et leur rapport), sa calibration (la
pente de la régression de Mincer--Zarnowitz proche de 1, la \omterm{def:rs:from-signal-to-forecast:curve}{courbe de calibration}), sa décroissance (chapitre~13), sa stabilité (chapitre~13), et sa valeur incrémentale par rapport aux
autres prévisions de la société (chapitre~12). La loi indique alors quel portefeuille attendre, et l'écart entre la loi et le backtest n'a rien de mystérieux~: il
est le bruit de l'IC, les contraintes et les coûts (chapitre~27), chacun mesurable. Michaud, Esch et Michaud (2020) ont soutenu que les applications de
la loi de Grinold à la conception de portefeuilles, en ajoutant des titres ou des facteurs ou en négociant plus souvent, sont souvent peu fiables parce qu'elles ignorent l'erreur d'estimation et
les contraintes de la pratique~; la décomposition ci-dessus est le contrôle qui garde un gestionnaire honnête.

\section{Fiches de prédicteur}

\begin{predictorcard}{Composite calibré}\label{pred:rs:from-signal-to-forecast:calibrated}
\sfield{Définition} $\alpha_i = \widehat{\mathrm{IC}}\,\omega_i z_i$, $z$ le composite du chapitre~14 standardisé par date, $\widehat{\mathrm{IC}}$ sa
moyenne glissante de l'IC, $\omega_i$ la volatilité résiduelle du modèle de risque sur l'horizon.
\sfield{Entrées et horodatages} Le composite et le modèle de risque à la date de décision~; l'IC estimé sur les dates antérieures seulement.
\sfield{Justification} \cref{prop:rs:from-signal-to-forecast:scaling}.
\sfield{Horizon et demi-vie} Celle du composite~; sur $h$ périodes, la somme de ses IC à une seule période.
\sfield{Normalisation} Déjà en unités de rendement~; vérifié chaque mois par la pente de Mincer--Zarnowitz.
\sfield{Modes de défaillance} Un IC estimé dans l'échantillon (in-sample)~; des volatilités issues d'un modèle de risque périmé~; des queues extrapolées à partir d'un ajustement linéaire.
\sfield{Sources} Grinold (1989)~; \texttt{firm.forecast}~; \texttt{rs\_forecast.calibration}.
\end{predictorcard}

\section{Tutoriel~: pourquoi le ratio d'information n'est pas ce que la loi promettait}\label{tut:rs:from-signal-to-forecast:tut}

\begin{tutorial}
\textbf{Objectif.} Calibrer un score artificiel en rendements espérés de trois façons, vérifier la loi avec des résidus gaussiens, puis mesurer l'écart à la promesse sur les
résidus du marché et sous contrainte long-only. \textbf{État final~:}
les \cref{fig:rs:from-signal-to-forecast:calibration,fig:rs:from-signal-to-forecast:shortfall}~; le tableau du tracking error.
\begin{tutsteps}
\item \textbf{\omterm{def:rs:from-signal-to-forecast:curve}{Régression isotone}} par regroupement des violateurs adjacents.
\omcode{code/firm/forecast/firm_forecast.py}{42}{62}{Regroupement des violateurs adjacents.}\label{lst:rs:from-signal-to-forecast:pav}
\item \textbf{La loi et ses corrections.}
\omcode{code/firm/forecast/firm_forecast.py}{73}{91}{Largeur effective, coefficient de transfert et formulations de la loi.}\label{lst:rs:from-signal-to-forecast:law}
\item \textbf{Exécuter} \texttt{calibration()}, \texttt{law()}, \texttt{tracking\_error\_grid()} et \texttt{fig\_forecast.py}.
\end{tutsteps}
\textbf{Pour aller plus loin.} Faire dépendre l'IC de la volatilité du titre (plus élevé pour les petits titres volatils) et observer que la règle de mise à l'échelle est mal calibrée~;
ajouter une contrainte de neutralité sectorielle et mesurer son \omterm{def:rs:from-signal-to-forecast:breadth}{coefficient de transfert}.
\end{tutorial}

\section{Construction~: la boîte à outils des prévisions}\label{bld:rs:from-signal-to-forecast:forecast}

\begin{build}
\bfield{Rôle} Tout composite qui parvient à un portefeuille arrive sous la forme d'un rendement espéré, avec sa calibration vérifiée et son ratio d'information espéré
décomposé.
\bfield{Interface} \texttt{scale\_rule(z, ic, vol)}, \texttt{binned}, \texttt{isotonic}, \texttt{mincer\_zarnowitz}, \texttt{effective\_breadth},
\texttt{transfer\_coefficient(w, alpha, vol)}, \texttt{law\_ir(ic, breadth, tc)}, \texttt{qian\_hua\_ir}, \texttt{ding\_martin\_ir},
\texttt{information\_ratio}.
\bfield{Règles} Les IC servant à la mise à l'échelle sont estimés avant les dates auxquelles ils s'appliquent~; les prévisions sont exprimées en unités de rendement par horizon~; chaque backtest présente la
promesse de la loi à côté de son ratio d'information réalisé.
\bfield{Tests d'acceptation} \texttt{code/firm/forecast/tests/}~: la règle de mise à l'échelle et les intervalles à la main~; \omterm{def:rs:from-signal-to-forecast:curve}{régression isotone} sur une suite donnée à la main, sur des entrées non
triées et sur une droite bruitée~; Mincer--Zarnowitz sur une pente artificielle~; \omterm{def:rs:from-signal-to-forecast:breadth}{largeur effective} aux extrêmes~; les formulations à la main~; un \omterm{def:rs:from-signal-to-forecast:breadth}{coefficient de transfert} égal à 1
pour le portefeuille sans contrainte et inférieur à 1 lorsque les positions courtes sont retirées.
\bfield{Pour aller plus loin} Calibration par régime~; IC par tranche de volatilité~; la forme de Ding et Martin avec un IC conditionnel estimé.
\end{build}

\begin{omsources}
\item R.~C.~Grinold, «~The fundamental law of active management~», \emph{Journal of Portfolio Management} 15(3), 1989.
\item R.~Clarke, H.~de Silva et S.~Thorley, «~Portfolio constraints and the fundamental law of active management~», \emph{Financial Analysts Journal}
58(5), 2002.
\item F.~Zhang, X.~Wang et H.~Cao, «~Turnover-adjusted information ratio~», arXiv:2105.10306, 2021 (la loi et ses formes de Qian--Hua et de Ding--Martin).
\item R.~O.~Michaud, D.~N.~Esch et R.~O.~Michaud, «~Estimation error and the fundamental law of active management~», \emph{Journal of Investing}
29(4), 2020.
\item M.~Ayer, H.~D.~Brunk, G.~M.~Ewing, W.~T.~Reid et E.~Silverman, «~An empirical distribution function for sampling with incomplete information~»,
\emph{Annals of Mathematical Statistics} 26(4), 1955.
\end{omsources}

\section{Exercices}

\begin{exercise}[$\star$]\label{exo:rs:from-signal-to-forecast:1}
Un score de $-1.5$, un IC de 0.04 et une volatilité résiduelle mensuelle de 10~\%~: quel est le \omterm{def:rs:anatomy-of-a-predictor:residual}{rendement résiduel} attendu pour le mois~? Et avec un score de
$+2$ et une volatilité de 5~\%~?
\end{exercise}

\begin{exercise}[$\star$]\label{exo:rs:from-signal-to-forecast:2}
Calculer le ratio d'information que promet la loi pour un IC de 0.03 sur 500 titres rééquilibrés chaque mois, puis sur 50 titres rééquilibrés chaque semaine (52 fois par an).
\end{exercise}

\begin{exercise}[$\star$]\label{exo:rs:from-signal-to-forecast:3}
Ajuster à la main une \omterm{def:rs:from-signal-to-forecast:curve}{régression isotone} aux valeurs 2, 1, 3, 3, 2, 4 pour des scores croissants.
\end{exercise}

\begin{exercise}[$\star\star$]\label{exo:rs:from-signal-to-forecast:4}
Calculer la \omterm{def:rs:from-signal-to-forecast:breadth}{largeur effective} de cent paris dont la corrélation moyenne deux à deux vaut 0.05, puis déterminer la corrélation qui la diviserait par deux.
\end{exercise}

\begin{exercise}[$\star\star$]\label{exo:rs:from-signal-to-forecast:5}
Un IC mensuel a pour moyenne 0.04 et pour écart type 0.05. Calculer le ratio d'information annuel prédit par la forme de Qian--Hua et par la loi avec 704 titres, et
expliquer pourquoi elles diffèrent.
\end{exercise}

\begin{exercise}[$\star\star$]\label{exo:rs:from-signal-to-forecast:6}
Les IC à un jour d'un signal valent 0.02, 0.01, puis 0 après le deuxième jour, avec une volatilité résiduelle quotidienne de 2~\%. Calculer le rendement attendu qu'implique un score de $+1$
sur un jour et sur un mois.
\end{exercise}

\begin{exercise}[$\star\star\star$]\label{exo:rs:from-signal-to-forecast:7}
\emph{Programmation.} Avec \texttt{rs\_forecast.portfolio}, mesurer le \omterm{def:rs:from-signal-to-forecast:breadth}{coefficient de transfert} et le ratio d'information du portefeuille long only à une tracking
error de 1~\% par mois lorsque l'indice de référence est concentré (poids proportionnels à une taille lognormale). Pourquoi cela change-t-il~?
\end{exercise}

\begin{exercise}[$\star\star\star$]\label{exo:rs:from-signal-to-forecast:8}
\emph{Trouver la faille.} «~Notre IC vaut 0.03 sur 3\,000 titres rééquilibrés chaque jour, donc notre ratio d'information devrait valoir $0.03\sqrt{252 \times 3000} = 26$.~»
\end{exercise}

\section{Problème~: pourquoi notre ratio d'information vaut-il les deux tiers de celui de la loi~?}

\begin{problem}[{Problème du week-end --- un écart, décomposé}]\label{pb:rs:from-signal-to-forecast:1}
La prévision de référence du chapitre~: IC 0.05, 704 titres, dix ans de données mensuelles, sur \texttt{firm.synthmkt}.

\medskip\noindent\textbf{Partie I --- Calibration.}
\begin{enumerate}
\item Quel rendement attendu un score de $+2$ implique-t-il pour un titre de volatilité médiane~?
\item Quel IC est estimé sur les cinq premières années~?
\item Donner la pente de la régression de Mincer--Zarnowitz et le $R^2$ de la prévision mise à l'échelle, ainsi que la pente sur le score brut.
\item Pourquoi le $R^2$ est-il si faible, et pourquoi cela n'a-t-il pas d'importance~?
\item Quand préférer la calibration isotone à la calibration linéaire~?
\end{enumerate}

\medskip\noindent\textbf{Partie II --- La loi.}
\begin{enumerate}[resume]
\item Que promet la loi~?
\item Que réalise le portefeuille avec des résidus gaussiens, et est-ce cohérent avec la loi~?
\item Quelles sont la moyenne et l'écart type de l'IC avec les résidus de marché et avec des résidus gaussiens~?
\item Quelle \omterm{def:rs:from-signal-to-forecast:breadth}{largeur effective} impliquent-elles~?
\end{enumerate}

\medskip\noindent\textbf{Partie III --- La contrainte.}
\begin{enumerate}[resume]
\item Quel \omterm{def:rs:from-signal-to-forecast:breadth}{coefficient de transfert} atteint le portefeuille long only à une tracking error demandée de 4~\% par mois, et quelle part des positions est écrêtée~?
\item Que prédit la loi généralisée, et quelle valeur est réalisée~?
\item Comment le \omterm{def:rs:from-signal-to-forecast:breadth}{coefficient de transfert} et le ratio d'information varient-ils avec la tracking error demandée~?
\item Énoncer le \emph{résultat nommé}~: la décomposition de l'écart en \omterm{def:rs:from-signal-to-forecast:breadth}{largeur effective}, \omterm{def:rs:from-signal-to-forecast:breadth}{coefficient de transfert} et bruit de la contrainte.
\end{enumerate}

\medskip\noindent\textbf{Partie IV --- Au-delà.}
\begin{enumerate}[resume]
\item Quel effet auraient les coûts sur la dernière barre~?
\item Comment un indice de référence pondéré par la capitalisation modifierait-il le \omterm{def:rs:from-signal-to-forecast:breadth}{coefficient de transfert}~?
\item Pourquoi ajouter des titres peut-il réduire le ratio d'information en pratique~?
\item Que doit publier la firme pour chaque stratégie~?
\item Laquelle des trois étapes la recherche peut-elle améliorer, et laquelle seule la construction de portefeuille peut améliorer~?
\item Quel est le danger d'utiliser la loi pour dimensionner une nouvelle stratégie~?
\item En une phrase~: que mesure la loi~?
\end{enumerate}
\end{problem}

\section{Questions d'entretien}

\begin{interviewq}[$\star$ \iqroles{chercheur}]\label{iq:rs:from-signal-to-forecast:1}
Énoncez la \omterm{def:rs:from-signal-to-forecast:law}{loi fondamentale de la gestion active} et ses hypothèses.
\end{interviewq}

\begin{interviewq}[$\star\star$ \iqroles{chercheur}]\label{iq:rs:from-signal-to-forecast:2}
Comment transformez-vous un score alpha en rendement attendu~?
\end{interviewq}

\begin{interviewq}[$\star\star$ \iqroles{chercheur, trader}]\label{iq:rs:from-signal-to-forecast:3}
Le ratio d'information de votre backtest vaut la moitié de ce que prédit la loi. Où regardez-vous~?
\end{interviewq}

\begin{interviewq}[$\star\star$ \iqroles{chercheur}]\label{iq:rs:from-signal-to-forecast:4}
Qu'est-ce que le \omterm{def:rs:from-signal-to-forecast:breadth}{coefficient de transfert}, et qu'est-ce qui le réduit~?
\end{interviewq}

\begin{interviewq}[$\star\star$ \iqroles{chercheur, ingénieur ML}]\label{iq:rs:from-signal-to-forecast:5}
Comment vérifieriez-vous que les prévisions de rendement d'un modèle sont calibrées~?
\end{interviewq}

\begin{interviewq}[$\star\star\star$ \iqroles{chercheur}]\label{iq:rs:from-signal-to-forecast:6}
Dériver la loi généralisée $\mathrm{IR} = \mathrm{TC}\cdot\mathrm{IC}\sqrt{\mathrm{BR}}$ et dire quelle hypothèse ne tient pas en premier en pratique.
\end{interviewq}
